% !TEX TS-program = lualatex

\DocumentMetadata{
  lang = en,
  pdfversion = 2.0,
  pdfstandard = ua-2,
  tagging = on,
  tagging-setup = {
    math/setup = {mathml-SE},
    table/header-rows = 1
  }
}

\documentclass[11pt]{article}

\usepackage{amsmath}
\usepackage{geometry}
\geometry{
  letterpaper,
  textwidth=6.5in,
  textheight=9in,
  top=0.85in,
  bottom=1in,
  left=0.75in,
  right=0.75in,
}
\usepackage{hyperref}

\newtheorem{theorem}{Theorem}
\newtheorem{acknowledgement}[theorem]{Acknowledgement}
\newtheorem{algorithm}[theorem]{Algorithm}
\newtheorem{axiom}[theorem]{Axiom}
\newtheorem{case}[theorem]{Case}
\newtheorem{claim}[theorem]{Claim}
\newtheorem{conclusion}[theorem]{Conclusion}
\newtheorem{condition}[theorem]{Condition}
\newtheorem{conjecture}[theorem]{Conjecture}
\newtheorem{corollary}[theorem]{Corollary}
\newtheorem{criterion}[theorem]{Criterion}
\newtheorem{definition}[theorem]{Definition}
\newtheorem{example}[theorem]{Example}
\newtheorem{exercise}[theorem]{Exercise}
\newtheorem{lemma}[theorem]{Lemma}
\newtheorem{notation}[theorem]{Notation}
\newtheorem{problem}[theorem]{Problem}
\newtheorem{proposition}[theorem]{Proposition}
\newtheorem{remark}[theorem]{Remark}
\newtheorem{solution}[theorem]{Solution}
\newtheorem{summary}[theorem]{Summary}
\newenvironment{proof}[1][Proof]{\textbf{#1.} }{\ \rule{0.5em}{0.5em}}

\tagpdfsetup{
  math/alt/use
}

\hypersetup{
  pdftitle={Assignment 1},
  pdfauthor={Blankenau},
  pdfsubject={Fall, 2026},
  hidelinks,
}

\begin{document}
\begin{center}
{\LARGE Assignment 1}
\end{center}

\noindent Blankenau

\noindent Economics 905, Fall, 2026

\medskip
\noindent \textit{Instructions.} You may work in groups and may compare
answers prior to submission. Joint submissions are allowed and encouraged.
All work should be legible and concise. A due date will be announced in
class.

\begin{enumerate}
\item Parts a-g follow your notes but try to work this without reference to
that for practice. Consider a one-period economy populated by a
representative agent and a representative firm. The agent is endowed with
one unit of time which is allocated between working, $n_{s},$ and leisure $%
\ell $. The agent makes these choices to maximize utility given by 
\begin{equation*}
u\left( c,\ell \right) =\frac{c^{1-\gamma }-1}{1-\gamma }+\frac{\ell
^{1-\gamma }-1}{1-\gamma }
\end{equation*}%
where $c$ is consumption and $\gamma >0$ measures the degree of curvature in
the utility function with respect to consumption and leisure. For each unit
of time worked, the agent is paid a wage $\omega $. The agent is also
endowed with $k_{s}$ units of capital which she rents inelastically in the
capital market. Each unit of capital rented pays $r$ and then fully
depreciates. The firm combines capital, $k_{d},$ and labor, $n_{d},$ to
produce a final good according to 
\begin{equation}
y\left( k,n\right) =zk_{d}^{\alpha }n_{d}^{1-\alpha },\text{ }0<\alpha <1
\label{out}
\end{equation}%
All markets are competitive.

\begin{enumerate}
\item Find the first order condition(s) for the agent's problem.

\item Find the first order conditions for the firm's problem.

\item Find the market clearing conditions.

\item Define an equilibrium.

\item Arrive at one equation and one unknown that could be used to solve for 
$\ell .$

\item Discuss how to find the other items of interest.

\item State and solve the social planner's problem. You will not be able to
get an explicit solution so arrive at one equation and 1 unknown. \ Will $%
\ell $ be the same as in the competitive equilibrium?

\item Suppose instead that there are two goods and preferences are given by 
\begin{equation*}
u\left( c,\ell \right) =\frac{c_{1}^{1-\gamma }-1}{1-\gamma }+\frac{%
c_{2}^{1-\gamma }-1}{1-\gamma }+\frac{\ell ^{1-\gamma }-1}{1-\gamma }
\end{equation*}%
where $c_{1}$ and $c_{2}$ are consumption of good 1 and good two. There is a
representative competitive firm producing each of the two goods according to 
\begin{eqnarray*}
y_{1} &=&z_{1}k_{1}^{\alpha }n_{1}^{1-\alpha } \\
y_{2} &=&z_{2}k_{2}^{\alpha }n_{2}^{1-\alpha }
\end{eqnarray*}%
where $n_{j}$ are $k_{j}$ are capital and labor hired by the firm producing
good $j$. Note that in equilibrium $k_{1}+k_{2}=k_{s}$ and $%
n_{1}+n_{2}=n_{s}.$ Write down the social planner's problem.

\item Find a system of three equations and three unknowns that could be used
to solve for the equilibrium values of $n_{1},n_{2},$ and $k_{1}.$

\item For the case where $\gamma =1$ find $n_{1},$ $n_{2},$ $\ell ,$ $k_{1}$
and $k_{2}.$
\end{enumerate}

\item Consider the following representative agent model. The consumer has
preferences given by 
\begin{equation*}
u\left( c,\ell \right) =\ln c+\ell
\end{equation*}%
where $c$ and $\ell $ are consumption and leisure$.$ The consumer is endowed
with 1 unit of time and $k_{0}$ units of capital. The representative firm
has a technology for producing consumption goods given by $y=zk^{\gamma
}n^{1-\gamma }$ where $z$ is total factor productivity, $k$ is the capital
input, $n$ is the labor input and $0<\gamma <1.$ Let $\omega $ denote the
wage rate and $r$ the rental rate on capital.

\begin{enumerate}
\item Define an equilibrium in this economy.

\item Solve for the competitive equilibrium prices and quantities.

\item Now suppose $y=zk+zn.$ Solve for the competitive equilibrium prices
and quantities. Here you will need to be careful about the requirement that $%
\ell \leq 1.$\pagebreak
\end{enumerate}
\end{enumerate}

\begin{enumerate}
\item 
\begin{enumerate}
\item The agent's problem is to maximize%
\begin{equation*}
u\left( c,\ell \right) =\frac{c^{1-\gamma }-1}{1-\gamma }+\frac{\ell
^{1-\gamma }-1}{1-\gamma }
\end{equation*}%
subject to 
\begin{equation*}
c\leq \omega \left( 1-\ell \right) +rk_{s}
\end{equation*}%
We know that this constraint will hold with equality so we can plug it in
above. It is also fine to set this up as a Lagrangian. Here we get 
\begin{equation*}
u\left( c,\ell \right) =\frac{\left[ \omega \left( 1-\ell \right) +rk_{s}%
\right] ^{1-\gamma }-1}{1-\gamma }+\frac{\ell ^{1-\gamma }-1}{1-\gamma }
\end{equation*}%
so that the FOC is 
\begin{equation}
-\omega \left[ \omega \left( 1-\ell \right) +rk_{s}\right] ^{-\gamma }+\ell
^{-\gamma }=0  \label{ut}
\end{equation}

\item The FOCs for the firm's problem are 
\begin{eqnarray}
r &=&\alpha zk_{d}^{\alpha -1}n_{d}^{1-\alpha }  \label{r} \\
\omega &=&\left( 1-\alpha \right) zk_{d}^{\alpha }n_{d}^{-\alpha }
\label{w2}
\end{eqnarray}

\begin{enumerate}
\item goods market clearing: 
\begin{equation}
y=c  \label{m1}
\end{equation}%
$.$

\item labor market clearing: 
\begin{equation}
n_{d}=n_{s}=\left( 1-\ell \right)  \label{m2}
\end{equation}%
$.$

\item capital market clearing:%
\begin{equation}
k_{s}=k_{d}  \label{m3}
\end{equation}%
$.$
\end{enumerate}

\item \textit{Definition:}\ A competitive equilibrium in this economy is a
set of quantities $\left\{ c,\ell ,n,k\right\} $ and prices $\left\{ \omega
,r\right\} $ which satisfy the following properties.

\begin{enumerate}
\item Each consumer chooses $c$ and $\ell $ optimally given $\omega $ and $%
r; $ i.e. equation (\ref{ut}) holds$.$

\item The representative firm chooses $n$ and $k$ optimally given $\omega $
and $r.$ i.e. equations (\ref{r}) and (\ref{w2}) hold.

\item Markets clear; i.e. equations (\ref{m1}), (\ref{m2}), and (\ref{m3})
hold.
\end{enumerate}

\item Equations (\ref{r}), (\ref{w2}), (\ref{m1}), (\ref{m2}), and (\ref{m3}%
) into equation (\ref{ut}) gives 
\begin{equation*}
-\left( 1-\alpha \right) zk_{d}^{\alpha }n_{d}^{-\alpha }\left[ \left(
1-\alpha \right) zk_{s}^{\alpha }\left( 1-\ell \right) ^{1-\alpha }+\alpha
zk_{s}^{\alpha }\left( 1-\ell \right) ^{1-\alpha }\right] ^{-\gamma }+\ell
^{-\gamma }=0
\end{equation*}%
or%
\begin{equation*}
\left( 1-\alpha \right) zk_{s}^{\alpha }\left( 1-\ell \right) ^{-\alpha } 
\left[ zk_{s}^{\alpha }\left( 1-\ell \right) ^{1-\alpha }\right] ^{-\gamma
}=\ell ^{-\gamma }
\end{equation*}%
or%
\begin{equation}
\left( 1-\alpha \right) z^{1-\gamma }k_{s}^{\alpha \left( 1-\gamma \right)
}\left( 1-\ell \right) ^{\alpha \gamma -\gamma -\alpha }=\ell ^{-\gamma }
\label{ell}
\end{equation}

\begin{enumerate}
\item use equation (\ref{ell}) to find $\ell .$

\item use this in equation (\ref{m2}) to find $n_{d}.$

\item use (\ref{m3}) to find $k_{d}.$ Then $y$, $r$, and $\omega $ can be
found from equations (\ref{out}), (\ref{r}), and (\ref{w2}).

\item Finally equation (\ref{m1}) gives $c.$
\end{enumerate}

\item From your notes, the social planner solves:%
\begin{equation*}
u\left( c,\ell \right) =\frac{\left( zk^{\alpha }\left( 1-\ell \right)
^{1-\alpha }\right) ^{1-\gamma }-1}{1-\gamma }+\frac{\ell ^{1-\gamma }-1}{%
1-\gamma }
\end{equation*}%
Setting the FOC equal to zero gives:%
\begin{equation}
\left[ zk_{o}^{\alpha }\left( 1-\ell \right) ^{1-\alpha }\right] ^{-\gamma
}\left( 1-\alpha \right) zk_{o}^{\alpha }\left( 1-\ell \right) ^{-\alpha
}=\ell ^{-\gamma }  \label{sp}
\end{equation}%
Comparing equations (\ref{ell})\ and (\ref{sp}), these are clearly the same.

\item You can do this in a number of ways but I plugged in constraints to get%
\begin{equation*}
\underset{k_{1},n_{1},n_{2}}{\max }\frac{\left( z_{1}k_{1}^{\alpha
}n_{1}^{1-\alpha }\right) ^{1-\gamma }-1}{1-\gamma }+\frac{\left(
z_{2}\left( k_{s}-k_{1}\right) ^{\alpha }n_{2}^{1-\alpha }\right) ^{1-\gamma
}-1}{1-\gamma }+\frac{\left( 1-n_{1}-n_{2}\right) ^{1-\gamma }-1}{1-\gamma }
\end{equation*}

\item The FOCs are 
\begin{eqnarray*}
\left( 1-\alpha \right) \left( z_{1}k_{1}^{\alpha }\right) ^{1-\gamma
}n_{1}^{\left( 1-\alpha \right) \left( 1-\gamma \right) -1} &=&\left(
1-n_{1}-n_{2}\right) ^{-\gamma } \\
\left( 1-\alpha \right) \left( z_{2}\left( k_{s}-k_{1}\right) ^{\alpha
}\right) ^{1-\gamma }n_{2}^{\left( 1-\alpha \right) \left( 1-\gamma \right)
-1} &=&\left( 1-n_{1}-n_{2}\right) ^{-\gamma } \\
\left( z_{2}n_{2}^{1-\alpha }\right) ^{1-\gamma }\left( k_{s}-k_{1}\right)
^{\alpha \left( 1-\gamma \right) -1} &=&\left( z_{1}n_{1}^{1-\alpha }\right)
^{1-\gamma }k_{1}^{\alpha \left( 1-\gamma \right) -1}
\end{eqnarray*}%
Notice that the only unknowns are $n_{1},n_{2}$ and $k_{1}.$

\item In the special case these are%
\begin{equation*}
\left( 1-\alpha \right) n_{1}^{-1}=\left( 1-n_{1}-n_{2}\right) ^{-1}
\end{equation*}%
\begin{equation*}
\left( 1-\alpha \right) n_{2}^{-1}=\left( 1-n_{1}-n_{2}\right) ^{-1}
\end{equation*}%
\begin{equation*}
k_{1}^{-1}=\left( k_{s}-k_{1}\right) ^{-1}
\end{equation*}%
From the first two%
\begin{equation*}
n_{1}=n_{2}
\end{equation*}%
so that the first gives 
\begin{equation*}
\left( 1-\alpha \right) \left( 1-2n_{2}\right) =n_{2}
\end{equation*}%
Thus%
\begin{eqnarray*}
n_{2} &=&\frac{1-\alpha }{3-2\alpha } \\
n_{1} &=&\frac{1-\alpha }{3-2\alpha } \\
\ell &=&\frac{1}{3-2\alpha }
\end{eqnarray*}%
From $k_{1}^{-1}=\left( k_{s}-k_{1}\right) ^{-1},$ clearly $k_{1}=\frac{k_{s}%
}{2}.$
\end{enumerate}

\item A competitive equilibrium is a set of quantities $\left\{ c,\ell
,n,k\right\} $ and prices $\left\{ \omega ,r\right\} $ such that

\begin{enumerate}
\item 
\begin{itemize}
\item each consumer chooses $c$ and $\ell $ optimally given $\omega ,r$.

\item the representative firm chooses $k$ and $n$ optimally given $\omega $
and $r.$ That is the firm solves%
\begin{equation*}
\underset{k,n}{\max }zk^{\gamma }n^{1-\gamma }-\omega n-rk
\end{equation*}

\item Markets clear. That is:%
\begin{eqnarray*}
k &=&k_{0} \\
n+\ell &=&1 \\
c &=&y
\end{eqnarray*}
\end{itemize}

\item The consumer solves 
\begin{equation*}
\underset{\ell }{\max }\ln c+\ell
\end{equation*}%
st%
\begin{eqnarray*}
c &\leq &rk+\omega \left( 1-\ell \right) \\
k &\leq &k_{0} \\
0 &\leq &\ell \leq 1
\end{eqnarray*}%
The first two constraints will hold with equality in equilibrium we can put
them into the objective function to get 
\begin{equation*}
\underset{\ell }{\max }\ln \left( rk+\omega \left( 1-\ell \right) \right)
+\ell
\end{equation*}%
The FOC is%
\begin{equation}
\frac{\omega }{rk+\omega \left( 1-\ell \right) }=1  \label{foc}
\end{equation}%
The firm solves 
\begin{equation*}
\underset{k,n}{\max }f(c,\ell )-\omega n-rk
\end{equation*}%
\ so optimization on the part of the firm gives 
\begin{eqnarray}
r &=&z\gamma \left( \frac{n}{k}\right) ^{1-\gamma }  \label{wages} \\
\omega &=&z\left( 1-\gamma \right) \left( \frac{k}{n}\right) ^{\gamma } 
\notag
\end{eqnarray}%
Putting (\ref{wages})into (\ref{foc}) gives 
\begin{equation*}
\frac{z\left( 1-\gamma \right) \left( \frac{k}{n}\right) ^{\gamma }}{z\gamma
\left( \frac{n}{k}\right) ^{1-\gamma }k_{0}+z\left( 1-\gamma \right) \left( 
\frac{k}{n}\right) ^{\gamma }\left( 1-\ell \right) }=1
\end{equation*}%
in equilibrium $k_{0}=k$ and $\left( 1-\ell \right) =n$ so 
\begin{equation*}
\frac{z\left( 1-\gamma \right) \left( \frac{k}{n}\right) ^{\gamma }}{z\gamma
\left( \frac{n}{k}\right) ^{1-\gamma }k+z\left( 1-\gamma \right) \left( 
\frac{k}{n}\right) ^{\gamma }n}=1
\end{equation*}%
simplify to get 
\begin{equation*}
\frac{\left( 1-\gamma \right) }{n\gamma +\left( 1-\gamma \right) n}=1
\end{equation*}%
and solve to get 
\begin{equation*}
n=1-\gamma .
\end{equation*}%
so 
\begin{eqnarray*}
\ell &=&\gamma \\
k &=&k_{0} \\
c &=&y=zk^{\gamma }n^{1-\gamma } \\
\omega &=&z\left( 1-\gamma \right) \left( \frac{k}{n}\right) ^{\gamma } \\
r &=&z\gamma \left( \frac{n}{k}\right) ^{1-\gamma }
\end{eqnarray*}%
There are a couple of things to notice. First, we have solved for every item
that the definition calls for. Second, in listing $\ell $ through $r$ we
have done it in such a way that the expressions include no terms that have
not previously been defined in terms of the model's parameters. Third, we
have one constraint that we have not yet dealt with, $0\leq \ell \leq 1.$
However since $0<\gamma <1$ this is satisfied. \ Finally, we are sloppy here
in that we ignore second or conditions. (Are these solutions maximums or
minimums?) You can always ignore second order conditions when not explicitly
asked for. However, be aware that they are required to actually be rigorous.

\item Solution: \ Optimization on the part of the firm gives 
\begin{eqnarray}
r &=&z  \label{wages2} \\
\omega &=&z  \notag
\end{eqnarray}%
Equation (\ref{wages2}) into (\ref{foc}) gives 
\begin{equation*}
\frac{z}{zk_{0}+z\left( 1-\ell \right) }=1
\end{equation*}%
and in equilibrium $k_{0}=k$ and $\left( 1-\ell \right) =n$. Solving then
gives $n=1-k$ or $\ell =k$. Now our requirement $0\leq \ell \leq 1$ has not
necessarily been satisfied. Since $k>0$ the left hand side is satisfied.
However, $k$ can exceed 1. \ Thus our solution is 
\begin{equation*}
\ell =\left\{ 
\begin{array}{cc}
k & \text{if }k\leq 1 \\ 
1 & \text{otherwise}%
\end{array}%
\right. .
\end{equation*}
\end{enumerate}
\end{enumerate}
\end{document}
